\section{Historical inverse PINN study}
This is Yash's contribution, using the historical pelt series supplied by Advay. It
extends the separately preserved synthetic recovery experiment; it is not a complete
combined group report. The Stan compilation contains annual Hudson's Bay Company
hare and lynx pelt returns \cite{carpenter2018}. We train on 1900--1915 (16 pairs) and
evaluate once on frozen forecasts for 1916--1920 (five pairs). Counts are thousands of
pelts, treated as abundance indices rather than censuses. Annual totals are approximated
as point samples. No observation scale factor is fitted. Alternative source compilations
disagree; all results here are conditional on the Stan series.

\subsection{Model and protocol}
With $t=\mathrm{Year}-1900$, the shared dimensional model is
\begin{align*}
\dot H&=rH(1-H/K)-aHP/(b+H), & \dot P&=eaHP/(b+H)-mP.
\end{align*}
All six positive coefficients are inferred. Bounds below are exploratory numerical
choices, not measured biological ranges. No coexistence constraint is imposed and $e$
is not constrained below one because these are effective coefficients in pelt units.
The starting vector is $(r,K,a,b,e,m)=(0.8,150,2,30,0.6,0.6)$. Yash chose this protocol
with assistance; agreement by the other members is not claimed. The local least-squares
benchmark uses the identical config, data split, scaling and solver.
\begin{center}\begin{tabular}{lrr}\toprule
Coefficient & Bounds & Selected primary estimate\\\midrule
$r$ & 0.05--3 & 0.60546 \\
$K$ & 20--500 & 497.6 \\
$a$ & 0.05--10 & 9.2605 \\
$b$ & 0.5--300 & 299.98 \\
$e$ & 0.02--5 & 1.0966 \\
$m$ & 0.05--3 & 0.99388 \\
\bottomrule\end{tabular}\end{center}

\subsection{Network, loss and optimization}
The inverse PINN \cite{raissi2019} uses three hidden layers of 48 tanh units, float64,
softplus-positive normalized outputs, and sigmoid-bounded coefficients. Species scales
are training maxima $(s_H,s_P)=(77.4,59.4)$; network time is $2t/15-1$.
For $u_\theta=(H_\theta/s_H,P_\theta/s_P)$, automatic differentiation evaluates
$R_\theta=\partial_tu_\theta-f(u_\theta\odot s;p)/s$. The loss is
\[
L=\operatorname{mean}(u_\theta-y/s)^2+
10\operatorname{mean}R_\theta^2+
10\operatorname{mean}(u_\theta(0)-y_0/s)^2.
\]
The initial observation is $y_0=(30,4)$. There are 161 collocation times in $[0,15]$.
Optimization uses 500 data/IC warmup steps, up to 5000 joint Adam steps at learning rate
0.002 (gradient norm clipped at 10), then up to 600 L-BFGS iterations with strong-Wolfe
line search. Training loss is logged every 100 Adam steps; plateau stopping uses
15 checks and minimum improvement $10^{-8}$. The best recorded weighted training
objective determines the saved state. Seeds are 17, 29 and 43, with one CPU thread and
deterministic Torch operations. Each variant selects its own lowest-loss seed; objectives
from different observation models are not used to select the primary analysis.
\begin{figure}[htbp]\centering
\ReportGraphEPS[width=\linewidth]{yash_historical_fit.eps}
\caption{Training observations, learned network and re-integrated ODE. A small local
physics residual does not guarantee identical long-horizon ODE and network trajectories.}
\end{figure}

\subsection{Frozen forecasts and matched benchmark}
The selected primary seed is 17. Forecasts use the learned coefficients in the shared
RK45 solver with rtol $10^{-8}$, atol $10^{-10}$ and maximum step 0.1 year. Integration
starts from $(30,4)$ in 1900 and proceeds through 1920, with no restart from a test state.
The neural network is never extrapolated for forecasting. Its starting-state discrepancy
is $(0.10060,-0.00723)$ thousand pelts; the ODE still starts at the exact observation.
The primary ODE training normalized MSE is 0.00981. The matched primary least-squares
fit uses six starts (RNG seed 91), bounded \texttt{least\_squares}, up to 500 function
evaluations and ftol=xtol=gtol=$10^{-10}$. Each start and its convergence status is retained.
This is an independent benchmark computed for this contribution; Advay's own historical
baseline is not represented as completed. PINN and least squares differ in soft versus
exact IC enforcement and in the PINN's added physics penalty.
\begin{figure}[htbp]\centering
\ReportGraphEPS[width=\linewidth]{yash_historical_forecast.eps}
\caption{Frozen forecasts. Thin pale lines are the three primary PINN seeds, not confidence
bands. Shading marks held-out years; observations were used for scoring only after all
forecast and seed choices were frozen.}
\end{figure}
\begin{center}\small\begin{tabular}{lrrrr}\toprule
& \multicolumn{2}{c}{RMSE} & \multicolumn{2}{c}{MAE}\\
Method & Hare & Lynx & Hare & Lynx\\\midrule
PINN fitted IC & 3.687 & 4.697 & 3.057 & 4.530 \\
PINN log errors & 18.351 & 10.026 & 14.289 & 8.972 \\
PINN primary & 2.847 & 4.317 & 2.590 & 4.046 \\
PINN wider bounds & 2.976 & 4.053 & 2.721 & 3.779 \\
PINN physics 100 & 4.519 & 3.978 & 4.107 & 3.496 \\
LS primary & 3.072 & 3.043 & 2.586 & 2.559 \\
LS log errors & 6.874 & 1.596 & 6.626 & 1.381 \\
LS fitted IC & 2.246 & 3.236 & 2.083 & 2.713 \\
\bottomrule\end{tabular}\end{center}
All scores are thousands of pelts on the same five held-out years. The primary PINN has
slightly lower hare RMSE than primary least squares, but higher lynx RMSE. Hare MAE is
essentially tied. These results do not establish an overall PINN advantage.
\begin{figure}[htbp]\centering
\ReportGraphEPS[width=0.95\linewidth]{yash_historical_loss.eps}
\caption{Selected primary training loss components. L-BFGS abscissae count closure
evaluations after Adam; they are not directly comparable computational steps.}
\end{figure}

\subsection{Sensitivity and practical identifiability}
All sensitivities were specified before any test scoring and are reported regardless
of outcome. The main three seeds were repeated with log observation residuals, and
separately with trainable initial states in $[0.5,1.5]$ times the 1900 observations. The
first data point stays in the data loss. The fitted-IC PINN's ODE starts at approximately
$(34.7477,3.7841)$; the network and ODE initial states are separately recorded. The benchmark
repeats both observation/IC sensitivities with six starts each. Log residuals worsen hare
forecast error for both methods, especially the PINN. None of these scores changes the
predeclared primary choice.

To check arbitrary box limits, each lower bound was halved and each upper bound doubled,
and the main PINN was repeated for all three seeds. A further seed-17 run increased physics
weight from 10 to 100. In total there are 13 PINN runs and 18 least-squares starts.
The wider-box fits move $K$ from about 498 to about 986 and $b$ from about 300 to about 600,
close to the expanded limits; $a$ also shifts substantially. Similar results across seeds
therefore do not demonstrate uniquely recovered biological parameters. The wider-box
forecast changes are smaller than these coefficient shifts.
\begin{figure}[htbp]\centering
\ReportGraphEPS[width=\linewidth]{yash_historical_bounds.eps}
\caption{Primary versus wider-box parameter estimates over three seeds. Several effective
coefficients track the numerical bounds; narrow seed spread is not parameter certainty.}
\end{figure}
\begin{figure}[htbp]\centering
\ReportGraphEPS[width=\linewidth]{yash_historical_sensitivity.eps}
\caption{All selected PINN variants. The log-error fit shows appreciable model/ODE forecast
sensitivity. Variants were retained rather than chosen using held-out errors.}
\end{figure}

\clearpage
The local trajectory sensitivity matrix uses central differences with respect to log
parameters (step $10^{-4}$), normalized using the declared observation loss. Its selected
primary condition number is approximately 317.5. This is a local diagnostic, not
a proof of structural identifiability or a confidence interval. For feasible coexistence,
$H^*=mb/(ea-m)$, $K_H=b+2H^*$ and $\sigma=K/K_H$; the selected primary fit gives
$\sigma=1.363$. This fitted ODE lies on the cycling side of Joel's threshold, conditional
on the model and protocol. It does not prove that the historical ecological system has
a stable limit cycle. All seeds' values and feasibility flags are exported to the team.

\subsection{Reproducibility and limitations}
\path{configs/historical.json} records every primary choice; the two sensitivity
configs record the expanded box and physics weight. Run manifests hash sources, config,
training data, forecasts and diagnostics. A forecast-freeze record is written before the
test observations are read. \path{scripts/run_historical_pinn.py} accepts only the
training table; \path{scripts/score_historical.py} handles final scoring. Existing
synthetic recovery results are preserved as a separate controlled validation.

Six effective coefficients are fitted from roughly two noisy cycles; boundary hits,
source disagreement, annual season totals, unknown catchability and the small test window
limit interpretation. PINN physics is a soft constraint, and the log-loss variant exposes
an appreciable network/ODE discrepancy. Seed spread is optimization sensitivity only.
Forecast skill on five points cannot establish general superiority, causal mechanisms
or absolute animal populations. This report and notebook concern Yash's component;
the final group report and submission still require integration.
